\subsection{NORMED ALGEBRAS OVER A VALUED FIELD}

DEFINITION 9. If A is an algebra over a non-discrete valued field K, a norm \(p(x)\) on A (A being considered as a vector space over K) is said to be compatible with the algebra structure of A if the topology it defines is compatible with the ring structure of A. An algebra over K, endowed with the structure defined by a norm compatible with the algebra structure, is called a normed algebra.

If A is a normed algebra over K, and if \(p(\boldsymbol{x})\) is the norm on A, the bilinear mapping \((\boldsymbol{x}, \boldsymbol{y}) \to \boldsymbol{x}\boldsymbol{y}\) of A × A into A is continuous, by hypothesis; hence by Theorem 1 of no. 5 there exists a real number a \textgreater{} 0 such that \(p(\boldsymbol{x}\boldsymbol{y}) \leqslant a.p(\boldsymbol{x})p(\boldsymbol{y})\) . Replacing \(p(\boldsymbol{x})\) by the equivalent norm \(a.p(\boldsymbol{x})\) , we may therefore always assume that the norm \(||\boldsymbol{x}||\) on a normed algebra A is such that

\[
\left| \left| x y \right| \right| \leqslant \left| \left| x \right| \right|. \left| \left| y \right| \right|.\tag{5}
\]

It follows from (5) by induction that, for each integer \(n > 0\) , we have

\[
\left| \left| x ^ {n} \right| \right| \leqslant \left| \left| x \right| \right| ^ {n}.\tag{6}
\]

Examples. 1) Let K be a valued division ring and let \(K'\) be a subfield of the centre of K such that the trace on \(K'\) of the absolute value \(|x|\) of K is not the improper absolute value on \(K'\) . Then K, with \(|x|\) as norm, is a normed algebra over \(K'\) .

\begin{enumerate}
\def\labelenumi{\arabic{enumi})}
\setcounter{enumi}{1}
\item
  Let K be a non-discrete valued field and let \(\mathbf{M}_{n}(\mathbf{K})\) be the ring of square matrices of order n over K. Regarded as a vector space over K, \(\mathbf{M}_{n}(\mathbf{K})\) is isomorphic to \(K^{n^{2}}\) . If for each \(X = (x_{ij}) \in \mathbf{M}_{n}(\mathbf{K})\) we define \(\|X\| = \sup_{i,j} |x_{ij}|\) , then \(\|X\|\) is a norm on \(\mathbf{M}_{n}(\mathbf{K})\) , and the topology defined by this norm is the product topology on \(K^{n^{2}}\) (Proposition 10); from this it follows (because of the continuity of polynomials in any number of variables over K) that this norm is compatible with the K-algebra structure of \(\mathbf{M}_{n}(\mathbf{K})\) .
\item
  The set \(\mathcal{B}(\mathbf{X};\mathbf{K})\) of all functions \(f\) on a set \(\mathbf{X}\) with values in a non-discrete valued field \(\mathbf{K}\) , such that \(x\to |f(x)|\) is bounded on \(\mathbf{X}\) , is an algebra over \(\mathbf{K}\) ; the norm \(\| f\| = \sup_{x\in \mathbf{X}}|f(x)|\) is compatible with the ring structure of \(\mathcal{B}(\mathbf{X};\mathbf{K})\) , because we have \(\| fg\| \leqslant \| f\| .\| g\|\) (cf.~Chapter X, § 1).
\end{enumerate}

Let \(\mathfrak{a}\) be a closed two-sided ideal in a normed algebra \(\mathbf{A}\) . If in the quotient algebra \(\mathbf{A} / \mathfrak{a}\) we put \(||\dot{\mathbf{x}}|| = \inf_{\mathbf{x}\in \dot{\mathbf{x}}}||\mathbf{x}||\) , we get a norm on \(\mathbf{A} / \mathfrak{a}\) which defines the topology which is the quotient by \(\mathfrak{a}\) of the topology of \(\mathbf{A}\) (Proposition 9); since this quotient topology is compatible with the quotient ring structure of \(\mathbf{A} / \mathfrak{a}\) (Chapter III, § 6, no. 4) it follows that the quotient algebra \(\mathbf{A} / \mathfrak{a}\) , with the norm \(||\mathbf{x}||\) , is a normed algebra.

Likewise, if \((\mathbf{A}_i)_{1 \leqslant i \leqslant n}\) is a family of \(n\) normed algebras over a valued field \(\mathbf{K}\) , and if in the product algebra \(\mathbf{A} = \prod_{i=1}^{n} \mathbf{A}_i\) we put \(||\boldsymbol{x}|| = \sup_{i} ||\boldsymbol{x}_i||\) , where \(\boldsymbol{x} = (\boldsymbol{x}_i)\) , we have a norm on \(\mathbf{A}\) which defines the product of the topologies of the \(\mathbf{A}_i\) (Proposition 10); since this topology is compatible with the ring structure of \(\mathbf{A}\) (Chapter III, § 6, no. 4), it follows that the product algebra \(\mathbf{A}\) , with the norm \(||\boldsymbol{x}||\) , is a normed algebra.

Let A be a normed algebra over a valued field K. The completion \(\hat{\mathbf{A}}\) of A (Chapter III, § 6, no. 5, Proposition 6) is also endowed with a \(\hat{\mathbf{K}}\) -vector space structure (no. 3, Proposition 8), and it is clear from the principle of extension of identities that we have \(t(xy) = (tx)y = x(ty)\) for all \(t \in \hat{\mathbf{K}}\) and all \(x, y \in \hat{\mathbf{A}}\) . Hence \(\hat{\mathbf{A}}\) is an algebra over \(\hat{\mathbf{K}}\) ; on the other hand (no. 3, Proposition 8), the norm on A extends by continuity to a norm which defines the topology of \(\hat{\mathbf{A}}\) , and therefore \(\hat{\mathbf{A}}\) , endowed with this norm, is a normed algebra over the field \(\hat{\mathbf{K}}\) .

If \((\pmb{x}_{\lambda})_{\lambda \in \mathbf{L}}\) and \((\pmb{y}_{\mu})_{\mu \in \mathbf{M}}\) are two absolutely summable families in a normed algebra \(\mathbf{A}\) , the family \((\pmb{x}_{\lambda}\pmb{y}_{\mu})_{(\lambda, \mu) \in \mathbf{L} \times \mathbf{M}}\) is absolutely summable, because \(||\pmb{x}_{\lambda}\pmb{y}_{\mu}|| \leqslant ||\pmb{x}_{\lambda}|| \cdot ||\pmb{y}_{\mu}||\) (Chapter IV, § 7, no. 3, Proposition 1); if in addition \(\mathbf{A}\) is complete, all three families are summable and we have

\[
\sum_ {(\lambda , \mu) \in \mathbf {L} \times \mathbf {M}} x _ {\lambda} y _ {\mu} = \left(\sum_ {\lambda \in \mathbf {L}} x _ {\lambda}\right) \left(\sum_ {\mu \in \mathbf {M}} y _ {\mu}\right)
\]

by the associativity of the sum on the left-hand side (Chapter III, § 5, no. 3, formula (2)).

If the normed algebra A has an identity element \(e \neq 0\) , the mapping \(t \rightarrow te\) is an isomorphism of the field structure of K onto that of the subfield Ke of A; this isomorphism is also an isomorphism of the topological field structure of K onto that of Ke (the topology of the latter being induced by that of A), for the restriction \(\|te\|\) of the norm of A to Ke is a norm equivalent to the absolute value

\[
| t | = \frac {\mathrm{I}}{| | e | |} | | t e | |.
\]

If \(||e|| = 1\) we have \(||te|| = |t|\) , and we can then identify the valued field \(K\) with the normed subfield \(Ke\) of \(A\) , and in particular we may denote the identity element of \(A\) by the symbol \(I\) .

In what follows we shall be concerned only with normed algebras which have an identity element e, and in which the norm satisfies the inequality (5); putting x = y = e in this inequality, it follows that \(\|e\| \geqslant 1\) .

PROPOSITION 12. If the series whose general term is \(z^n\) is convergent in A, then \(e - z\) is a unit of A and we have

\[
(e - z) ^ {- 1} = \sum_ {n = 0} ^ {\infty} z ^ {n}.\tag{7}
\]

Conversely, if \(||z|| < 1\) and if \(e - z\) is a unit in \(A\) , then the series whose general term is \(z^n\) is convergent and formula (7) is valid.

For each \(p > 0\) we have

\[
(e - z) \sum_ {n = 0} ^ {p} z ^ {n} = e - z ^ {p + 1}.\tag{8}
\]

If the series whose general term is \(z^{n}\) is convergent and if y is its sum, then \(z^{n}\) tends to o as \(n \to +\infty\) ; hence by passing to the limit in (8) we have \((e - z)y = e\) ; similarly we prove that \(y(e - z) = e\) , and hence \(y = (e - z)^{-1}\) (note that this part of the argument is valid in any topological ring which has an identity element).

Conversely, if \(||z|| < 1\) , then since \(||z^{p+1}|| \leqslant ||z||^{p+1}\) , it follows that \(z^{p+1}\) tends to o as \(p \to +\infty\) ; multiplying both sides (8) on the left by \((e - z)^{-1}\) and letting \(p\) tend to infinity, we see that the series whose general term is \(z^n\) converges and has \((e - z)^{-1}\) as its sum.

COROLLARY. Let A be a complete normed algebra. Then for each \(z \in \mathbf{A}\) such that \(||z|| < 1\) , \(e - z\) is a unit in A.

The series whose general term is \(z^n\) is absolutely convergent, since \(||z^n|| \leqslant ||z||^n\) for \(n > 0\) , and is therefore convergent, since \(A\) is complete (no. 6, Proposition 11).

PROPOSITION 13. Let G be the group of units of a complete normed algebra A. Then G is an open subset of A; the topology induced on G by the topology of A is compatible with the group structure of G; and G, endowed with this topology, is a complete group (with respect to each of its two uniformities).

The corollary to Proposition 12 shows that G contains a neighbourhood V of e in A; hence, for each \(x_{0} \in G\) , the elements of \(x_{0}V\) are units, and \(x_{0}V\) is a neighbourhood of \(x_{0}\) in A, since \(x \to x_{0}x\) is a homeomorphism of A onto itself ( \(x_{0}\) being a unit of A). Hence G is open in A.

To show that the topology induced on \(\mathbf{G}\) by the topology of \(\mathbf{A}\) is compatible with the group structure of \(\mathbf{G}\) , it is sufficient to show that the function \(x^{-1}\) is continuous on \(\mathbf{G}\) . Let \(x_0 \in \mathbf{G}\) , and for each \(x \in \mathbf{G}\) , write \(\boldsymbol{x}\) in the form \(x = x_0 (e + u)\) , so that \(u = x_0^{-1}(x - x_0)\) ; then \(||\boldsymbol{u}|| \leqslant ||x_0^{-1}|| \cdot ||\boldsymbol{x} - x_0||\) , and thus if \(||\boldsymbol{x} - x_0|| \leqslant 1 / ||x_0^{-1}||\) , we have \(||\boldsymbol{u}|| \leqslant 1\) , \(e + u = x_0^{-1}x\) is a unit, the series whose general term is \((-1)^n u^n\) is absolutely convergent, and

\[
\boldsymbol {x} ^ {- 1} = (\boldsymbol {e} + \boldsymbol {u}) ^ {- 1} \boldsymbol {x} _ {0} ^ {- 1} = \boldsymbol {x} _ {0} ^ {- 1} + \left(\sum_ {n = 1} ^ {\infty} (- 1) ^ {n} \boldsymbol {u} ^ {n}\right) \boldsymbol {x} _ {0} ^ {- 1},\tag{9}
\]

from which it follows that

\[
\begin{array}{r l} & {| | x _ {0} ^ {- 1} - x _ {0} ^ {- 1} | | \leqslant \left\| \sum_ {n = 0} ^ {\infty} (- \mathrm{I}) ^ {n} u ^ {n} \right\|. | | u | |. | | x _ {0} ^ {- 1} | |} \\ & {\quad \leqslant \left\| \sum_ {n = 0} ^ {\infty} (- \mathrm{I}) ^ {n} u ^ {n} \right\|. | | x _ {0} ^ {- 1} | | ^ {2}. | | x - x _ {0} | |.} \end{array}
\]

As \(x\) tends to \(x_0\) , \(||x - x_0||\) tends to \(o\) , and since

\[
\left\| \sum_ {n = 0} ^ {\infty} (- 1) ^ {n} \boldsymbol {u} ^ {n} \right\| \leqslant | | \boldsymbol {e} | | + \frac {| | \boldsymbol {u} | |}{1 - | | \boldsymbol {u} | |}
\]

remains bounded, it follows that \(x^{-1}\) tends to \(x_{0}^{-1}\) .

Finally, to show that the left uniformity of G is complete, let us show that every Cauchy filter F with respect to this uniformity is a Cauchy filter with respect to the additive uniformity of A and converges to a point of G. For each \(\varepsilon\) such that \(0 < \varepsilon < 1\) , there is a set \(\mathbf{M} \in \mathfrak{F}\) such that \(||x^{-1}y - e|| \leqslant \varepsilon\) for all \(x, y\) in M, i.e., such that

\[
\left| \left| \mathbf {y} - \mathbf {x} \right| \right| \leqslant \varepsilon \left| \left| \mathbf {x} \right| \right|.
\]

Let \(\pmb{a}\) be a point of M; for each \(x \in M\) we have \(||x - a|| \leqslant \varepsilon ||a||\) , and therefore \(||x|| \leqslant (1 + \varepsilon)||a||\) . On the other hand, there exists a set \(N \subset M\) belonging to \(\mathfrak{F}\) and such that \(||x^{-1}y - e|| \leqslant \varepsilon / (1 + \varepsilon)||a||\) for all \(x\) and \(y\) in N; it follows that \(||y - x|| \leqslant \varepsilon ||x|| / (1 + \varepsilon)||a|| \leqslant \varepsilon\) , which shows that \(\mathfrak{F}\) is a Cauchy filter with respect to the additive uniformity of A, and therefore converges to a point \(x_0\) , since A is complete. Since \(x_0\) is the limit of \(\mathfrak{F}\) , we have \(||x^{-1}x_0 - e|| \leqslant \varepsilon\) for all \(x \in M\) by the principle of extension of inequalities; since \(\varepsilon > 1\) , it follows that \(x^{-1}x_0\) is a unit in A; hence \(x_0\) is a unit, i.e., \(x_0 \in G\) .

PROPOSITION 14. In a complete valued division ring, the multiplicative group of non-zero elements is a complete group.

The proof is similar to that of Proposition 13; we have only to replace the norm of A by the absolute value of the division ring under consideration.

Note that we cannot apply Proposition 13 directly, because a (non-commutative) valued division ring is not necessarily an algebra over a non-discrete valued field (the restriction of the absolute value to the centre of the division ring might be improper).

Remark. Proposition 13 is not necessarily true for a non-complete normed algebra. For example, in the algebra \(\mathcal{C}(\mathbf{I};\mathbf{R})\) of all finite continuous real-valued functions on \(\mathbf{I} = [\mathrm{o},\mathbf{i}]\) {[}the norm being \(\| x\| = \sup_{t\in \mathbf{I}}|x(t)|\) , the subalgebra P of all polynomials in t (restricted to I) is not complete; if \(x(t)\) is any non-constant polynomial, then \(1 + \varepsilon x\) is arbitrarily close to the identity element i of P when \(\varepsilon\) is arbitrarily small, but \(1 + \varepsilon x\) is not a unit in P. However, if A is a non-complete normed algebra, G its group of units, \(\hat{A}\) the completion of A, then G is a subgroup of the group of units of \(\hat{A}\) , and therefore the topology induced on G by the topology of A is compatible with the group structure of G.

